% v1.1
% figures-lib.tex — shared figure macro library for paper + presentation.
% Spec: docs/superpowers/specs/2026-04-24-figures-lib-design.md
%
% Usage (in document preamble, AFTER \usepackage{tikz}):
%   % v1.1
% figures-lib.tex — shared figure macro library for paper + presentation.
% Spec: docs/superpowers/specs/2026-04-24-figures-lib-design.md
%
% Usage (in document preamble, AFTER \usepackage{tikz}):
%   % v1.1
% figures-lib.tex — shared figure macro library for paper + presentation.
% Spec: docs/superpowers/specs/2026-04-24-figures-lib-design.md
%
% Usage (in document preamble, AFTER \usepackage{tikz}):
%   % v1.1
% figures-lib.tex — shared figure macro library for paper + presentation.
% Spec: docs/superpowers/specs/2026-04-24-figures-lib-design.md
%
% Usage (in document preamble, AFTER \usepackage{tikz}):
%   \input{../figures-lib.tex}
%   \figstyle{paper}      % or \figstyle{slides}
%
% Rules:
%   • \input in preamble only — never inside tikzpicture / \fitwidth / \resizebox.
%   • Caller positions every primitive; primitives never pick their own coords.
%   • Style via figlib/* keys; never hardcode thick / \scriptsize at call sites.
%
% To extend: copy a primitive's section, document its bounding box, use figlib/*
% styles. Atoms = simple \newcommand. Pics = TikZ pic with keyval options.

\RequirePackage{etoolbox}
\usetikzlibrary{arrows.meta, calc, positioning, shapes.geometric}

% ─── Palette (idempotent: only defines if undefined) ────────────────────
\providecolor{uciGold}{HTML}{F0C040}
\providecolor{uciGoldPale}{HTML}{FAF0CC}
\providecolor{uciBlue}{HTML}{0064A4}
\providecolor{uciNavy}{HTML}{1B3D6F}
\providecolor{uciCream}{HTML}{FFF8F0}
\providecolor{uciAmber}{HTML}{C8830A}
\providecolor{uciSand}{HTML}{EED8A0}

% ─── Style switch ────────────────────────────────────────────────────────
\tikzset{
  figlib/line/.style={thick},
  figlib/thin/.style={semithick},
  figlib/font/.style={font=\scriptsize},
  figlib/tip/.style={>={Stealth[length=2mm]}},
}

\newcommand{\figstyle}[1]{%
  \ifstrequal{#1}{slides}{%
    \tikzset{%
      figlib/line/.style={thick},
      figlib/thin/.style={semithick},
      figlib/font/.style={font=\scriptsize},
      figlib/tip/.style={>={Stealth[length=2mm]}},
    }%
  }{%
    \tikzset{%
      figlib/line/.style={semithick},
      figlib/thin/.style={thin},
      figlib/font/.style={font=\footnotesize},
      figlib/tip/.style={>={Stealth[length=1.5mm]}},
    }%
  }%
}

% ============================================================
% Atoms — simple \newcommand, color via single optional arg.
% ============================================================

% \phasor[<color>]{<angle-deg>}{<magnitude>}
% Bbox: from (0,0) extending {magnitude} units at {angle}.
\newcommand{\phasor}[3][uciBlue]{%
  \draw[->, figlib/line, figlib/tip, draw=#1] (0,0) -- (#2:#3);
}

% \phasorL[<color>]{<angle-deg>}{<magnitude>}{<label>}
\newcommand{\phasorL}[4][uciBlue]{%
  \draw[->, figlib/line, figlib/tip, draw=#1] (0,0) -- (#2:#3);
  \node[figlib/font, color=#1] at (#2:{#3+0.28}) {#4};
}

% \phaseArc[<color>]{<from-deg>}{<to-deg>}{<radius>}{<label>}
% Bbox: arc of {radius} between the two angles.
\newcommand{\phaseArc}[5][uciBlue]{%
  \draw[figlib/thin, draw=#1]
    (#2:#4) arc[start angle=#2, end angle=#3, radius=#4];
  \node[figlib/font, color=#1]
    at ({(#2+#3)/2}:{#4*1.18}) {#5};
}

% \signal[<color>]{<from-coord>}{<to-coord>}
% from/to include parens, e.g. (A.east) or (1,2).
\newcommand{\signal}[3][uciNavy]{%
  \draw[->, figlib/line, figlib/tip, draw=#1] #2 -- #3;
}

% \signalL[<color>]{<from>}{<to>}{<label>}
\newcommand{\signalL}[4][uciNavy]{%
  \draw[->, figlib/line, figlib/tip, draw=#1] #2 -- #3
    node[midway, above, figlib/font, color=#1]{#4};
}

% ============================================================
% Pics — keyval options. All centered at the placement coord.
% ============================================================

% pic{block={label=..., width=..., height=...}}
% Bbox: width × height (cm, bare numbers). Default 1.0 × 0.7.
\tikzset{
  pics/block/.style={code={%
    \tikzset{block opts/.cd, label=, width=1.0, height=0.7, #1}%
    \pgfmathsetmacro{\figlibBw}{
      \pgfkeysvalueof{/tikz/block opts/width}/2}%
    \pgfmathsetmacro{\figlibBh}{
      \pgfkeysvalueof{/tikz/block opts/height}/2}%
    \draw[draw=uciNavy, fill=white, figlib/line]
      (-\figlibBw,-\figlibBh) rectangle (\figlibBw,\figlibBh);
    \node[figlib/font, color=uciNavy] at (0,0)
      {\pgfkeysvalueof{/tikz/block opts/label}};
  }},
  block opts/.cd,
    label/.initial=,
    width/.initial=1.0,
    height/.initial=0.7,
}

% pic{sum node}     — '+' in 0.5cm circle.
\tikzset{
  pics/sum node/.style={code={%
    \draw[draw=uciNavy, fill=white, figlib/line]
      (0,0) circle[radius=0.25cm];
    \node[figlib/font, color=uciNavy] at (0,0) {$+$};
  }},
}

% pic{mixer}        — '×' in 0.5cm circle.
\tikzset{
  pics/mixer/.style={code={%
    \draw[draw=uciNavy, fill=white, figlib/line]
      (0,0) circle[radius=0.25cm];
    \node[figlib/font, color=uciNavy] at (0,0) {$\times$};
  }},
}

% pic{branch node}  — 3pt filled dot.
\tikzset{
  pics/branch node/.style={code={%
    \fill[uciNavy] (0,0) circle[radius=1.5pt];
  }},
}

% pic{delay block={cycles=N, width=..., height=..., label=...}}
% N=0 → straight wire; N≥1 → that many sine cycles inside the box.
% Wire enters top mid-x, exits bottom mid-x. Label drawn LEFT of box.
% Bbox: width × height (cm, bare numbers). Default 0.55 × 0.9.
\tikzset{
  pics/delay block/.style={code={%
    \tikzset{delayblk opts/.cd,
      cycles=0, width=0.55, height=0.9, label=, #1}%
    \pgfmathsetmacro{\figlibDhalf}{
      \pgfkeysvalueof{/tikz/delayblk opts/height}/2}%
    \pgfmathsetmacro{\figlibDhalfW}{
      \pgfkeysvalueof{/tikz/delayblk opts/width}/2}%
    \pgfmathtruncatemacro{\figlibDcyc}{
      \pgfkeysvalueof{/tikz/delayblk opts/cycles}}%
    \draw[draw=uciNavy, fill=white, figlib/line]
      (-\figlibDhalfW,-\figlibDhalf)
      rectangle (\figlibDhalfW,\figlibDhalf);
    \node[figlib/font, color=uciNavy, anchor=east, inner sep=2pt]
      at (-\figlibDhalfW - 0.05, 0)
      {\pgfkeysvalueof{/tikz/delayblk opts/label}};
    \ifnum\figlibDcyc=0
      \draw[figlib/line, draw=uciBlue]
        (0,\figlibDhalf) -- (0,-\figlibDhalf);
    \else
      \pgfmathsetmacro{\figlibDmar}{0.05*
        \pgfkeysvalueof{/tikz/delayblk opts/height}}%
      \pgfmathsetmacro{\figlibDspan}{
        \pgfkeysvalueof{/tikz/delayblk opts/height}-2*\figlibDmar}%
      \pgfmathsetmacro{\figlibDstep}{\figlibDspan/(4*\figlibDcyc)}%
      \pgfmathsetmacro{\figlibDtop}{\figlibDhalf-\figlibDmar}%
      \draw[figlib/line, draw=uciBlue]
        (0,\figlibDhalf) -- (0,\figlibDtop)
        \foreach \figlibDk in {1,...,\figlibDcyc}{
          \pgfextra{%
            \pgfmathsetmacro{\figlibDya}{
              \figlibDtop-(4*\figlibDk-3)*\figlibDstep}%
            \pgfmathsetmacro{\figlibDyb}{
              \figlibDtop-(4*\figlibDk-2)*\figlibDstep}%
            \pgfmathsetmacro{\figlibDyc}{
              \figlibDtop-(4*\figlibDk-1)*\figlibDstep}%
            \pgfmathsetmacro{\figlibDyd}{
              \figlibDtop-(4*\figlibDk)*\figlibDstep}%
          }
          sin (0.13,\figlibDya) cos (0,\figlibDyb)
          sin (-0.13,\figlibDyc) cos (0,\figlibDyd)
        }
        -- (0,-\figlibDhalf);
    \fi
  }},
  delayblk opts/.cd,
    cycles/.initial=0,
    width/.initial=0.55,
    height/.initial=0.9,
    label/.initial=,
}

% pic{antenna array={N=4, spacing=1.0, labels=true, label prefix=ant\space}}
% First triangle at (0,0); subsequent at (k*spacing, 0). spacing in cm (bare).
% Bbox: x ∈ [−0.18, (N−1)*spacing+0.18], y ∈ [-0.4, 0.32] (incl labels).
\tikzset{
  pics/antenna array/.style={code={%
    \tikzset{antarr opts/.cd,
      N=4, spacing=1.0, labels=true, label prefix=ant\space, #1}%
    \pgfmathtruncatemacro{\figlibAN}{
      \pgfkeysvalueof{/tikz/antarr opts/N}-1}%
    \foreach \figlibAi in {0,...,\figlibAN}{%
      \pgfmathsetmacro{\figlibAx}{
        \figlibAi*\pgfkeysvalueof{/tikz/antarr opts/spacing}}%
      \draw[draw=uciNavy, fill=uciSand, figlib/line]
        (\figlibAx,0) -- (\figlibAx-0.18,0.32)
        -- (\figlibAx+0.18,0.32) -- cycle;
      \ifdefstring{\pgfkeysvalueof{/tikz/antarr opts/labels}}{true}{%
        \node[figlib/font, color=uciNavy, anchor=north]
          at (\figlibAx,-0.05)
          {\pgfkeysvalueof{/tikz/antarr opts/label prefix}\figlibAi};%
      }{}%
    }%
  }},
  antarr opts/.cd,
    N/.initial=4,
    spacing/.initial=1.0,
    labels/.initial=true,
    label prefix/.initial=ant\space,
}

% pic{wavefronts={N=3, length=2.4, angle=45, sep=0.7, arrow=true}}
% N parallel lines at {angle} from +x, length L, separated by sep along normal.
% First line passes through (0,0); subsequent shifted by sep along the normal.
% Optional propagation arrow drawn from (0,0) outward along -angle.
\tikzset{
  pics/wavefronts/.style={code={%
    \tikzset{wf opts/.cd,
      N=3, length=2.4, angle=45, sep=0.7, arrow=true, #1}%
    \pgfmathtruncatemacro{\figlibWN}{
      \pgfkeysvalueof{/tikz/wf opts/N}-1}%
    \pgfmathsetmacro{\figlibWnx}{
      cos(\pgfkeysvalueof{/tikz/wf opts/angle}+90)}%
    \pgfmathsetmacro{\figlibWny}{
      sin(\pgfkeysvalueof{/tikz/wf opts/angle}+90)}%
    \foreach \figlibWi in {0,...,\figlibWN}{%
      \pgfmathsetmacro{\figlibWdx}{
        \figlibWi*\pgfkeysvalueof{/tikz/wf opts/sep}*\figlibWnx}%
      \pgfmathsetmacro{\figlibWdy}{
        \figlibWi*\pgfkeysvalueof{/tikz/wf opts/sep}*\figlibWny}%
      \draw[figlib/line, draw=uciBlue!70]
        (\figlibWdx,\figlibWdy) --
          ++(\pgfkeysvalueof{/tikz/wf opts/angle}:%
             \pgfkeysvalueof{/tikz/wf opts/length});
    }%
    \ifdefstring{\pgfkeysvalueof{/tikz/wf opts/arrow}}{true}{%
      \pgfmathsetmacro{\figlibWpx}{
        0.5*\pgfkeysvalueof{/tikz/wf opts/length}*
          cos(\pgfkeysvalueof{/tikz/wf opts/angle})}%
      \pgfmathsetmacro{\figlibWpy}{
        0.5*\pgfkeysvalueof{/tikz/wf opts/length}*
          sin(\pgfkeysvalueof{/tikz/wf opts/angle})}%
      \draw[->, figlib/line, figlib/tip, draw=uciBlue]
        ($(\figlibWpx,\figlibWpy)+%
          (\pgfkeysvalueof{/tikz/wf opts/angle}+90:0.7)$)
          -- (\figlibWpx,\figlibWpy);
    }{}%
  }},
  wf opts/.cd,
    N/.initial=3,
    length/.initial=2.4,
    angle/.initial=45,
    sep/.initial=0.7,
    arrow/.initial=true,
}

% pic{beam lobe={angle=90, scale=1.0, sidelobes=true, color=uciAmber}}
% Lobe base at (0,0), main lobe pointing along {angle} from +x.
% Bbox: roughly 2*scale long, 1*scale wide.
\tikzset{
  pics/beam lobe/.style={code={%
    \tikzset{beam opts/.cd,
      angle=90, scale=1.0, sidelobes=true, color=uciAmber, #1}%
    \begin{scope}[shift={(0,0)},
                  rotate=\pgfkeysvalueof{/tikz/beam opts/angle}-90,
                  scale=\pgfkeysvalueof{/tikz/beam opts/scale}]
      \fill[\pgfkeysvalueof{/tikz/beam opts/color}!30,
            draw=\pgfkeysvalueof{/tikz/beam opts/color}, figlib/line]
        (0,0) .. controls (0.9,2.4) and (-0.9,2.4) .. (0,0);
      \ifdefstring{\pgfkeysvalueof{/tikz/beam opts/sidelobes}}{true}{%
        \fill[\pgfkeysvalueof{/tikz/beam opts/color}!30,
              draw=\pgfkeysvalueof{/tikz/beam opts/color}, figlib/line]
          (0,0) .. controls (0.9,0.7) and (0.35,0.7) .. (0,0);
        \fill[\pgfkeysvalueof{/tikz/beam opts/color}!30,
              draw=\pgfkeysvalueof{/tikz/beam opts/color}, figlib/line]
          (0,0) .. controls (-0.9,0.7) and (-0.35,0.7) .. (0,0);
      }{}%
    \end{scope}
  }},
  beam opts/.cd,
    angle/.initial=90,
    scale/.initial=1.0,
    sidelobes/.initial=true,
    color/.initial=uciAmber,
}

% pic{axes xy={xmin=-1, xmax=1, ymin=-1, ymax=1, xlabel=$x$, ylabel=$y$}}
% Origin at placement coord. Bbox: (xmax-xmin) × (ymax-ymin).
\tikzset{
  pics/axes xy/.style={code={%
    \tikzset{axes opts/.cd,
      xlabel=$x$, ylabel=$y$,
      xmin=-1, xmax=1, ymin=-1, ymax=1, #1}%
    \draw[->, figlib/thin, figlib/tip, draw=uciNavy]
      (\pgfkeysvalueof{/tikz/axes opts/xmin},0) --
      (\pgfkeysvalueof{/tikz/axes opts/xmax},0)
      node[right, figlib/font]{\pgfkeysvalueof{/tikz/axes opts/xlabel}};
    \draw[->, figlib/thin, figlib/tip, draw=uciNavy]
      (0,\pgfkeysvalueof{/tikz/axes opts/ymin}) --
      (0,\pgfkeysvalueof{/tikz/axes opts/ymax})
      node[above, figlib/font]{\pgfkeysvalueof{/tikz/axes opts/ylabel}};
  }},
  axes opts/.cd,
    xlabel/.initial=$x$,
    ylabel/.initial=$y$,
    xmin/.initial=-1, xmax/.initial=1,
    ymin/.initial=-1, ymax/.initial=1,
}

\endinput

%   \figstyle{paper}      % or \figstyle{slides}
%
% Rules:
%   • \input in preamble only — never inside tikzpicture / \fitwidth / \resizebox.
%   • Caller positions every primitive; primitives never pick their own coords.
%   • Style via figlib/* keys; never hardcode thick / \scriptsize at call sites.
%
% To extend: copy a primitive's section, document its bounding box, use figlib/*
% styles. Atoms = simple \newcommand. Pics = TikZ pic with keyval options.

\RequirePackage{etoolbox}
\usetikzlibrary{arrows.meta, calc, positioning, shapes.geometric}

% ─── Palette (idempotent: only defines if undefined) ────────────────────
\providecolor{uciGold}{HTML}{F0C040}
\providecolor{uciGoldPale}{HTML}{FAF0CC}
\providecolor{uciBlue}{HTML}{0064A4}
\providecolor{uciNavy}{HTML}{1B3D6F}
\providecolor{uciCream}{HTML}{FFF8F0}
\providecolor{uciAmber}{HTML}{C8830A}
\providecolor{uciSand}{HTML}{EED8A0}

% ─── Style switch ────────────────────────────────────────────────────────
\tikzset{
  figlib/line/.style={thick},
  figlib/thin/.style={semithick},
  figlib/font/.style={font=\scriptsize},
  figlib/tip/.style={>={Stealth[length=2mm]}},
}

\newcommand{\figstyle}[1]{%
  \ifstrequal{#1}{slides}{%
    \tikzset{%
      figlib/line/.style={thick},
      figlib/thin/.style={semithick},
      figlib/font/.style={font=\scriptsize},
      figlib/tip/.style={>={Stealth[length=2mm]}},
    }%
  }{%
    \tikzset{%
      figlib/line/.style={semithick},
      figlib/thin/.style={thin},
      figlib/font/.style={font=\footnotesize},
      figlib/tip/.style={>={Stealth[length=1.5mm]}},
    }%
  }%
}

% ============================================================
% Atoms — simple \newcommand, color via single optional arg.
% ============================================================

% \phasor[<color>]{<angle-deg>}{<magnitude>}
% Bbox: from (0,0) extending {magnitude} units at {angle}.
\newcommand{\phasor}[3][uciBlue]{%
  \draw[->, figlib/line, figlib/tip, draw=#1] (0,0) -- (#2:#3);
}

% \phasorL[<color>]{<angle-deg>}{<magnitude>}{<label>}
\newcommand{\phasorL}[4][uciBlue]{%
  \draw[->, figlib/line, figlib/tip, draw=#1] (0,0) -- (#2:#3);
  \node[figlib/font, color=#1] at (#2:{#3+0.28}) {#4};
}

% \phaseArc[<color>]{<from-deg>}{<to-deg>}{<radius>}{<label>}
% Bbox: arc of {radius} between the two angles.
\newcommand{\phaseArc}[5][uciBlue]{%
  \draw[figlib/thin, draw=#1]
    (#2:#4) arc[start angle=#2, end angle=#3, radius=#4];
  \node[figlib/font, color=#1]
    at ({(#2+#3)/2}:{#4*1.18}) {#5};
}

% \signal[<color>]{<from-coord>}{<to-coord>}
% from/to include parens, e.g. (A.east) or (1,2).
\newcommand{\signal}[3][uciNavy]{%
  \draw[->, figlib/line, figlib/tip, draw=#1] #2 -- #3;
}

% \signalL[<color>]{<from>}{<to>}{<label>}
\newcommand{\signalL}[4][uciNavy]{%
  \draw[->, figlib/line, figlib/tip, draw=#1] #2 -- #3
    node[midway, above, figlib/font, color=#1]{#4};
}

% ============================================================
% Pics — keyval options. All centered at the placement coord.
% ============================================================

% pic{block={label=..., width=..., height=...}}
% Bbox: width × height (cm, bare numbers). Default 1.0 × 0.7.
\tikzset{
  pics/block/.style={code={%
    \tikzset{block opts/.cd, label=, width=1.0, height=0.7, #1}%
    \pgfmathsetmacro{\figlibBw}{
      \pgfkeysvalueof{/tikz/block opts/width}/2}%
    \pgfmathsetmacro{\figlibBh}{
      \pgfkeysvalueof{/tikz/block opts/height}/2}%
    \draw[draw=uciNavy, fill=white, figlib/line]
      (-\figlibBw,-\figlibBh) rectangle (\figlibBw,\figlibBh);
    \node[figlib/font, color=uciNavy] at (0,0)
      {\pgfkeysvalueof{/tikz/block opts/label}};
  }},
  block opts/.cd,
    label/.initial=,
    width/.initial=1.0,
    height/.initial=0.7,
}

% pic{sum node}     — '+' in 0.5cm circle.
\tikzset{
  pics/sum node/.style={code={%
    \draw[draw=uciNavy, fill=white, figlib/line]
      (0,0) circle[radius=0.25cm];
    \node[figlib/font, color=uciNavy] at (0,0) {$+$};
  }},
}

% pic{mixer}        — '×' in 0.5cm circle.
\tikzset{
  pics/mixer/.style={code={%
    \draw[draw=uciNavy, fill=white, figlib/line]
      (0,0) circle[radius=0.25cm];
    \node[figlib/font, color=uciNavy] at (0,0) {$\times$};
  }},
}

% pic{branch node}  — 3pt filled dot.
\tikzset{
  pics/branch node/.style={code={%
    \fill[uciNavy] (0,0) circle[radius=1.5pt];
  }},
}

% pic{delay block={cycles=N, width=..., height=..., label=...}}
% N=0 → straight wire; N≥1 → that many sine cycles inside the box.
% Wire enters top mid-x, exits bottom mid-x. Label drawn LEFT of box.
% Bbox: width × height (cm, bare numbers). Default 0.55 × 0.9.
\tikzset{
  pics/delay block/.style={code={%
    \tikzset{delayblk opts/.cd,
      cycles=0, width=0.55, height=0.9, label=, #1}%
    \pgfmathsetmacro{\figlibDhalf}{
      \pgfkeysvalueof{/tikz/delayblk opts/height}/2}%
    \pgfmathsetmacro{\figlibDhalfW}{
      \pgfkeysvalueof{/tikz/delayblk opts/width}/2}%
    \pgfmathtruncatemacro{\figlibDcyc}{
      \pgfkeysvalueof{/tikz/delayblk opts/cycles}}%
    \draw[draw=uciNavy, fill=white, figlib/line]
      (-\figlibDhalfW,-\figlibDhalf)
      rectangle (\figlibDhalfW,\figlibDhalf);
    \node[figlib/font, color=uciNavy, anchor=east, inner sep=2pt]
      at (-\figlibDhalfW - 0.05, 0)
      {\pgfkeysvalueof{/tikz/delayblk opts/label}};
    \ifnum\figlibDcyc=0
      \draw[figlib/line, draw=uciBlue]
        (0,\figlibDhalf) -- (0,-\figlibDhalf);
    \else
      \pgfmathsetmacro{\figlibDmar}{0.05*
        \pgfkeysvalueof{/tikz/delayblk opts/height}}%
      \pgfmathsetmacro{\figlibDspan}{
        \pgfkeysvalueof{/tikz/delayblk opts/height}-2*\figlibDmar}%
      \pgfmathsetmacro{\figlibDstep}{\figlibDspan/(4*\figlibDcyc)}%
      \pgfmathsetmacro{\figlibDtop}{\figlibDhalf-\figlibDmar}%
      \draw[figlib/line, draw=uciBlue]
        (0,\figlibDhalf) -- (0,\figlibDtop)
        \foreach \figlibDk in {1,...,\figlibDcyc}{
          \pgfextra{%
            \pgfmathsetmacro{\figlibDya}{
              \figlibDtop-(4*\figlibDk-3)*\figlibDstep}%
            \pgfmathsetmacro{\figlibDyb}{
              \figlibDtop-(4*\figlibDk-2)*\figlibDstep}%
            \pgfmathsetmacro{\figlibDyc}{
              \figlibDtop-(4*\figlibDk-1)*\figlibDstep}%
            \pgfmathsetmacro{\figlibDyd}{
              \figlibDtop-(4*\figlibDk)*\figlibDstep}%
          }
          sin (0.13,\figlibDya) cos (0,\figlibDyb)
          sin (-0.13,\figlibDyc) cos (0,\figlibDyd)
        }
        -- (0,-\figlibDhalf);
    \fi
  }},
  delayblk opts/.cd,
    cycles/.initial=0,
    width/.initial=0.55,
    height/.initial=0.9,
    label/.initial=,
}

% pic{antenna array={N=4, spacing=1.0, labels=true, label prefix=ant\space}}
% First triangle at (0,0); subsequent at (k*spacing, 0). spacing in cm (bare).
% Bbox: x ∈ [−0.18, (N−1)*spacing+0.18], y ∈ [-0.4, 0.32] (incl labels).
\tikzset{
  pics/antenna array/.style={code={%
    \tikzset{antarr opts/.cd,
      N=4, spacing=1.0, labels=true, label prefix=ant\space, #1}%
    \pgfmathtruncatemacro{\figlibAN}{
      \pgfkeysvalueof{/tikz/antarr opts/N}-1}%
    \foreach \figlibAi in {0,...,\figlibAN}{%
      \pgfmathsetmacro{\figlibAx}{
        \figlibAi*\pgfkeysvalueof{/tikz/antarr opts/spacing}}%
      \draw[draw=uciNavy, fill=uciSand, figlib/line]
        (\figlibAx,0) -- (\figlibAx-0.18,0.32)
        -- (\figlibAx+0.18,0.32) -- cycle;
      \ifdefstring{\pgfkeysvalueof{/tikz/antarr opts/labels}}{true}{%
        \node[figlib/font, color=uciNavy, anchor=north]
          at (\figlibAx,-0.05)
          {\pgfkeysvalueof{/tikz/antarr opts/label prefix}\figlibAi};%
      }{}%
    }%
  }},
  antarr opts/.cd,
    N/.initial=4,
    spacing/.initial=1.0,
    labels/.initial=true,
    label prefix/.initial=ant\space,
}

% pic{wavefronts={N=3, length=2.4, angle=45, sep=0.7, arrow=true}}
% N parallel lines at {angle} from +x, length L, separated by sep along normal.
% First line passes through (0,0); subsequent shifted by sep along the normal.
% Optional propagation arrow drawn from (0,0) outward along -angle.
\tikzset{
  pics/wavefronts/.style={code={%
    \tikzset{wf opts/.cd,
      N=3, length=2.4, angle=45, sep=0.7, arrow=true, #1}%
    \pgfmathtruncatemacro{\figlibWN}{
      \pgfkeysvalueof{/tikz/wf opts/N}-1}%
    \pgfmathsetmacro{\figlibWnx}{
      cos(\pgfkeysvalueof{/tikz/wf opts/angle}+90)}%
    \pgfmathsetmacro{\figlibWny}{
      sin(\pgfkeysvalueof{/tikz/wf opts/angle}+90)}%
    \foreach \figlibWi in {0,...,\figlibWN}{%
      \pgfmathsetmacro{\figlibWdx}{
        \figlibWi*\pgfkeysvalueof{/tikz/wf opts/sep}*\figlibWnx}%
      \pgfmathsetmacro{\figlibWdy}{
        \figlibWi*\pgfkeysvalueof{/tikz/wf opts/sep}*\figlibWny}%
      \draw[figlib/line, draw=uciBlue!70]
        (\figlibWdx,\figlibWdy) --
          ++(\pgfkeysvalueof{/tikz/wf opts/angle}:%
             \pgfkeysvalueof{/tikz/wf opts/length});
    }%
    \ifdefstring{\pgfkeysvalueof{/tikz/wf opts/arrow}}{true}{%
      \pgfmathsetmacro{\figlibWpx}{
        0.5*\pgfkeysvalueof{/tikz/wf opts/length}*
          cos(\pgfkeysvalueof{/tikz/wf opts/angle})}%
      \pgfmathsetmacro{\figlibWpy}{
        0.5*\pgfkeysvalueof{/tikz/wf opts/length}*
          sin(\pgfkeysvalueof{/tikz/wf opts/angle})}%
      \draw[->, figlib/line, figlib/tip, draw=uciBlue]
        ($(\figlibWpx,\figlibWpy)+%
          (\pgfkeysvalueof{/tikz/wf opts/angle}+90:0.7)$)
          -- (\figlibWpx,\figlibWpy);
    }{}%
  }},
  wf opts/.cd,
    N/.initial=3,
    length/.initial=2.4,
    angle/.initial=45,
    sep/.initial=0.7,
    arrow/.initial=true,
}

% pic{beam lobe={angle=90, scale=1.0, sidelobes=true, color=uciAmber}}
% Lobe base at (0,0), main lobe pointing along {angle} from +x.
% Bbox: roughly 2*scale long, 1*scale wide.
\tikzset{
  pics/beam lobe/.style={code={%
    \tikzset{beam opts/.cd,
      angle=90, scale=1.0, sidelobes=true, color=uciAmber, #1}%
    \begin{scope}[shift={(0,0)},
                  rotate=\pgfkeysvalueof{/tikz/beam opts/angle}-90,
                  scale=\pgfkeysvalueof{/tikz/beam opts/scale}]
      \fill[\pgfkeysvalueof{/tikz/beam opts/color}!30,
            draw=\pgfkeysvalueof{/tikz/beam opts/color}, figlib/line]
        (0,0) .. controls (0.9,2.4) and (-0.9,2.4) .. (0,0);
      \ifdefstring{\pgfkeysvalueof{/tikz/beam opts/sidelobes}}{true}{%
        \fill[\pgfkeysvalueof{/tikz/beam opts/color}!30,
              draw=\pgfkeysvalueof{/tikz/beam opts/color}, figlib/line]
          (0,0) .. controls (0.9,0.7) and (0.35,0.7) .. (0,0);
        \fill[\pgfkeysvalueof{/tikz/beam opts/color}!30,
              draw=\pgfkeysvalueof{/tikz/beam opts/color}, figlib/line]
          (0,0) .. controls (-0.9,0.7) and (-0.35,0.7) .. (0,0);
      }{}%
    \end{scope}
  }},
  beam opts/.cd,
    angle/.initial=90,
    scale/.initial=1.0,
    sidelobes/.initial=true,
    color/.initial=uciAmber,
}

% pic{axes xy={xmin=-1, xmax=1, ymin=-1, ymax=1, xlabel=$x$, ylabel=$y$}}
% Origin at placement coord. Bbox: (xmax-xmin) × (ymax-ymin).
\tikzset{
  pics/axes xy/.style={code={%
    \tikzset{axes opts/.cd,
      xlabel=$x$, ylabel=$y$,
      xmin=-1, xmax=1, ymin=-1, ymax=1, #1}%
    \draw[->, figlib/thin, figlib/tip, draw=uciNavy]
      (\pgfkeysvalueof{/tikz/axes opts/xmin},0) --
      (\pgfkeysvalueof{/tikz/axes opts/xmax},0)
      node[right, figlib/font]{\pgfkeysvalueof{/tikz/axes opts/xlabel}};
    \draw[->, figlib/thin, figlib/tip, draw=uciNavy]
      (0,\pgfkeysvalueof{/tikz/axes opts/ymin}) --
      (0,\pgfkeysvalueof{/tikz/axes opts/ymax})
      node[above, figlib/font]{\pgfkeysvalueof{/tikz/axes opts/ylabel}};
  }},
  axes opts/.cd,
    xlabel/.initial=$x$,
    ylabel/.initial=$y$,
    xmin/.initial=-1, xmax/.initial=1,
    ymin/.initial=-1, ymax/.initial=1,
}

\endinput

%   \figstyle{paper}      % or \figstyle{slides}
%
% Rules:
%   • \input in preamble only — never inside tikzpicture / \fitwidth / \resizebox.
%   • Caller positions every primitive; primitives never pick their own coords.
%   • Style via figlib/* keys; never hardcode thick / \scriptsize at call sites.
%
% To extend: copy a primitive's section, document its bounding box, use figlib/*
% styles. Atoms = simple \newcommand. Pics = TikZ pic with keyval options.

\RequirePackage{etoolbox}
\usetikzlibrary{arrows.meta, calc, positioning, shapes.geometric}

% ─── Palette (idempotent: only defines if undefined) ────────────────────
\providecolor{uciGold}{HTML}{F0C040}
\providecolor{uciGoldPale}{HTML}{FAF0CC}
\providecolor{uciBlue}{HTML}{0064A4}
\providecolor{uciNavy}{HTML}{1B3D6F}
\providecolor{uciCream}{HTML}{FFF8F0}
\providecolor{uciAmber}{HTML}{C8830A}
\providecolor{uciSand}{HTML}{EED8A0}

% ─── Style switch ────────────────────────────────────────────────────────
\tikzset{
  figlib/line/.style={thick},
  figlib/thin/.style={semithick},
  figlib/font/.style={font=\scriptsize},
  figlib/tip/.style={>={Stealth[length=2mm]}},
}

\newcommand{\figstyle}[1]{%
  \ifstrequal{#1}{slides}{%
    \tikzset{%
      figlib/line/.style={thick},
      figlib/thin/.style={semithick},
      figlib/font/.style={font=\scriptsize},
      figlib/tip/.style={>={Stealth[length=2mm]}},
    }%
  }{%
    \tikzset{%
      figlib/line/.style={semithick},
      figlib/thin/.style={thin},
      figlib/font/.style={font=\footnotesize},
      figlib/tip/.style={>={Stealth[length=1.5mm]}},
    }%
  }%
}

% ============================================================
% Atoms — simple \newcommand, color via single optional arg.
% ============================================================

% \phasor[<color>]{<angle-deg>}{<magnitude>}
% Bbox: from (0,0) extending {magnitude} units at {angle}.
\newcommand{\phasor}[3][uciBlue]{%
  \draw[->, figlib/line, figlib/tip, draw=#1] (0,0) -- (#2:#3);
}

% \phasorL[<color>]{<angle-deg>}{<magnitude>}{<label>}
\newcommand{\phasorL}[4][uciBlue]{%
  \draw[->, figlib/line, figlib/tip, draw=#1] (0,0) -- (#2:#3);
  \node[figlib/font, color=#1] at (#2:{#3+0.28}) {#4};
}

% \phaseArc[<color>]{<from-deg>}{<to-deg>}{<radius>}{<label>}
% Bbox: arc of {radius} between the two angles.
\newcommand{\phaseArc}[5][uciBlue]{%
  \draw[figlib/thin, draw=#1]
    (#2:#4) arc[start angle=#2, end angle=#3, radius=#4];
  \node[figlib/font, color=#1]
    at ({(#2+#3)/2}:{#4*1.18}) {#5};
}

% \signal[<color>]{<from-coord>}{<to-coord>}
% from/to include parens, e.g. (A.east) or (1,2).
\newcommand{\signal}[3][uciNavy]{%
  \draw[->, figlib/line, figlib/tip, draw=#1] #2 -- #3;
}

% \signalL[<color>]{<from>}{<to>}{<label>}
\newcommand{\signalL}[4][uciNavy]{%
  \draw[->, figlib/line, figlib/tip, draw=#1] #2 -- #3
    node[midway, above, figlib/font, color=#1]{#4};
}

% ============================================================
% Pics — keyval options. All centered at the placement coord.
% ============================================================

% pic{block={label=..., width=..., height=...}}
% Bbox: width × height (cm, bare numbers). Default 1.0 × 0.7.
\tikzset{
  pics/block/.style={code={%
    \tikzset{block opts/.cd, label=, width=1.0, height=0.7, #1}%
    \pgfmathsetmacro{\figlibBw}{
      \pgfkeysvalueof{/tikz/block opts/width}/2}%
    \pgfmathsetmacro{\figlibBh}{
      \pgfkeysvalueof{/tikz/block opts/height}/2}%
    \draw[draw=uciNavy, fill=white, figlib/line]
      (-\figlibBw,-\figlibBh) rectangle (\figlibBw,\figlibBh);
    \node[figlib/font, color=uciNavy] at (0,0)
      {\pgfkeysvalueof{/tikz/block opts/label}};
  }},
  block opts/.cd,
    label/.initial=,
    width/.initial=1.0,
    height/.initial=0.7,
}

% pic{sum node}     — '+' in 0.5cm circle.
\tikzset{
  pics/sum node/.style={code={%
    \draw[draw=uciNavy, fill=white, figlib/line]
      (0,0) circle[radius=0.25cm];
    \node[figlib/font, color=uciNavy] at (0,0) {$+$};
  }},
}

% pic{mixer}        — '×' in 0.5cm circle.
\tikzset{
  pics/mixer/.style={code={%
    \draw[draw=uciNavy, fill=white, figlib/line]
      (0,0) circle[radius=0.25cm];
    \node[figlib/font, color=uciNavy] at (0,0) {$\times$};
  }},
}

% pic{branch node}  — 3pt filled dot.
\tikzset{
  pics/branch node/.style={code={%
    \fill[uciNavy] (0,0) circle[radius=1.5pt];
  }},
}

% pic{delay block={cycles=N, width=..., height=..., label=...}}
% N=0 → straight wire; N≥1 → that many sine cycles inside the box.
% Wire enters top mid-x, exits bottom mid-x. Label drawn LEFT of box.
% Bbox: width × height (cm, bare numbers). Default 0.55 × 0.9.
\tikzset{
  pics/delay block/.style={code={%
    \tikzset{delayblk opts/.cd,
      cycles=0, width=0.55, height=0.9, label=, #1}%
    \pgfmathsetmacro{\figlibDhalf}{
      \pgfkeysvalueof{/tikz/delayblk opts/height}/2}%
    \pgfmathsetmacro{\figlibDhalfW}{
      \pgfkeysvalueof{/tikz/delayblk opts/width}/2}%
    \pgfmathtruncatemacro{\figlibDcyc}{
      \pgfkeysvalueof{/tikz/delayblk opts/cycles}}%
    \draw[draw=uciNavy, fill=white, figlib/line]
      (-\figlibDhalfW,-\figlibDhalf)
      rectangle (\figlibDhalfW,\figlibDhalf);
    \node[figlib/font, color=uciNavy, anchor=east, inner sep=2pt]
      at (-\figlibDhalfW - 0.05, 0)
      {\pgfkeysvalueof{/tikz/delayblk opts/label}};
    \ifnum\figlibDcyc=0
      \draw[figlib/line, draw=uciBlue]
        (0,\figlibDhalf) -- (0,-\figlibDhalf);
    \else
      \pgfmathsetmacro{\figlibDmar}{0.05*
        \pgfkeysvalueof{/tikz/delayblk opts/height}}%
      \pgfmathsetmacro{\figlibDspan}{
        \pgfkeysvalueof{/tikz/delayblk opts/height}-2*\figlibDmar}%
      \pgfmathsetmacro{\figlibDstep}{\figlibDspan/(4*\figlibDcyc)}%
      \pgfmathsetmacro{\figlibDtop}{\figlibDhalf-\figlibDmar}%
      \draw[figlib/line, draw=uciBlue]
        (0,\figlibDhalf) -- (0,\figlibDtop)
        \foreach \figlibDk in {1,...,\figlibDcyc}{
          \pgfextra{%
            \pgfmathsetmacro{\figlibDya}{
              \figlibDtop-(4*\figlibDk-3)*\figlibDstep}%
            \pgfmathsetmacro{\figlibDyb}{
              \figlibDtop-(4*\figlibDk-2)*\figlibDstep}%
            \pgfmathsetmacro{\figlibDyc}{
              \figlibDtop-(4*\figlibDk-1)*\figlibDstep}%
            \pgfmathsetmacro{\figlibDyd}{
              \figlibDtop-(4*\figlibDk)*\figlibDstep}%
          }
          sin (0.13,\figlibDya) cos (0,\figlibDyb)
          sin (-0.13,\figlibDyc) cos (0,\figlibDyd)
        }
        -- (0,-\figlibDhalf);
    \fi
  }},
  delayblk opts/.cd,
    cycles/.initial=0,
    width/.initial=0.55,
    height/.initial=0.9,
    label/.initial=,
}

% pic{antenna array={N=4, spacing=1.0, labels=true, label prefix=ant\space}}
% First triangle at (0,0); subsequent at (k*spacing, 0). spacing in cm (bare).
% Bbox: x ∈ [−0.18, (N−1)*spacing+0.18], y ∈ [-0.4, 0.32] (incl labels).
\tikzset{
  pics/antenna array/.style={code={%
    \tikzset{antarr opts/.cd,
      N=4, spacing=1.0, labels=true, label prefix=ant\space, #1}%
    \pgfmathtruncatemacro{\figlibAN}{
      \pgfkeysvalueof{/tikz/antarr opts/N}-1}%
    \foreach \figlibAi in {0,...,\figlibAN}{%
      \pgfmathsetmacro{\figlibAx}{
        \figlibAi*\pgfkeysvalueof{/tikz/antarr opts/spacing}}%
      \draw[draw=uciNavy, fill=uciSand, figlib/line]
        (\figlibAx,0) -- (\figlibAx-0.18,0.32)
        -- (\figlibAx+0.18,0.32) -- cycle;
      \ifdefstring{\pgfkeysvalueof{/tikz/antarr opts/labels}}{true}{%
        \node[figlib/font, color=uciNavy, anchor=north]
          at (\figlibAx,-0.05)
          {\pgfkeysvalueof{/tikz/antarr opts/label prefix}\figlibAi};%
      }{}%
    }%
  }},
  antarr opts/.cd,
    N/.initial=4,
    spacing/.initial=1.0,
    labels/.initial=true,
    label prefix/.initial=ant\space,
}

% pic{wavefronts={N=3, length=2.4, angle=45, sep=0.7, arrow=true}}
% N parallel lines at {angle} from +x, length L, separated by sep along normal.
% First line passes through (0,0); subsequent shifted by sep along the normal.
% Optional propagation arrow drawn from (0,0) outward along -angle.
\tikzset{
  pics/wavefronts/.style={code={%
    \tikzset{wf opts/.cd,
      N=3, length=2.4, angle=45, sep=0.7, arrow=true, #1}%
    \pgfmathtruncatemacro{\figlibWN}{
      \pgfkeysvalueof{/tikz/wf opts/N}-1}%
    \pgfmathsetmacro{\figlibWnx}{
      cos(\pgfkeysvalueof{/tikz/wf opts/angle}+90)}%
    \pgfmathsetmacro{\figlibWny}{
      sin(\pgfkeysvalueof{/tikz/wf opts/angle}+90)}%
    \foreach \figlibWi in {0,...,\figlibWN}{%
      \pgfmathsetmacro{\figlibWdx}{
        \figlibWi*\pgfkeysvalueof{/tikz/wf opts/sep}*\figlibWnx}%
      \pgfmathsetmacro{\figlibWdy}{
        \figlibWi*\pgfkeysvalueof{/tikz/wf opts/sep}*\figlibWny}%
      \draw[figlib/line, draw=uciBlue!70]
        (\figlibWdx,\figlibWdy) --
          ++(\pgfkeysvalueof{/tikz/wf opts/angle}:%
             \pgfkeysvalueof{/tikz/wf opts/length});
    }%
    \ifdefstring{\pgfkeysvalueof{/tikz/wf opts/arrow}}{true}{%
      \pgfmathsetmacro{\figlibWpx}{
        0.5*\pgfkeysvalueof{/tikz/wf opts/length}*
          cos(\pgfkeysvalueof{/tikz/wf opts/angle})}%
      \pgfmathsetmacro{\figlibWpy}{
        0.5*\pgfkeysvalueof{/tikz/wf opts/length}*
          sin(\pgfkeysvalueof{/tikz/wf opts/angle})}%
      \draw[->, figlib/line, figlib/tip, draw=uciBlue]
        ($(\figlibWpx,\figlibWpy)+%
          (\pgfkeysvalueof{/tikz/wf opts/angle}+90:0.7)$)
          -- (\figlibWpx,\figlibWpy);
    }{}%
  }},
  wf opts/.cd,
    N/.initial=3,
    length/.initial=2.4,
    angle/.initial=45,
    sep/.initial=0.7,
    arrow/.initial=true,
}

% pic{beam lobe={angle=90, scale=1.0, sidelobes=true, color=uciAmber}}
% Lobe base at (0,0), main lobe pointing along {angle} from +x.
% Bbox: roughly 2*scale long, 1*scale wide.
\tikzset{
  pics/beam lobe/.style={code={%
    \tikzset{beam opts/.cd,
      angle=90, scale=1.0, sidelobes=true, color=uciAmber, #1}%
    \begin{scope}[shift={(0,0)},
                  rotate=\pgfkeysvalueof{/tikz/beam opts/angle}-90,
                  scale=\pgfkeysvalueof{/tikz/beam opts/scale}]
      \fill[\pgfkeysvalueof{/tikz/beam opts/color}!30,
            draw=\pgfkeysvalueof{/tikz/beam opts/color}, figlib/line]
        (0,0) .. controls (0.9,2.4) and (-0.9,2.4) .. (0,0);
      \ifdefstring{\pgfkeysvalueof{/tikz/beam opts/sidelobes}}{true}{%
        \fill[\pgfkeysvalueof{/tikz/beam opts/color}!30,
              draw=\pgfkeysvalueof{/tikz/beam opts/color}, figlib/line]
          (0,0) .. controls (0.9,0.7) and (0.35,0.7) .. (0,0);
        \fill[\pgfkeysvalueof{/tikz/beam opts/color}!30,
              draw=\pgfkeysvalueof{/tikz/beam opts/color}, figlib/line]
          (0,0) .. controls (-0.9,0.7) and (-0.35,0.7) .. (0,0);
      }{}%
    \end{scope}
  }},
  beam opts/.cd,
    angle/.initial=90,
    scale/.initial=1.0,
    sidelobes/.initial=true,
    color/.initial=uciAmber,
}

% pic{axes xy={xmin=-1, xmax=1, ymin=-1, ymax=1, xlabel=$x$, ylabel=$y$}}
% Origin at placement coord. Bbox: (xmax-xmin) × (ymax-ymin).
\tikzset{
  pics/axes xy/.style={code={%
    \tikzset{axes opts/.cd,
      xlabel=$x$, ylabel=$y$,
      xmin=-1, xmax=1, ymin=-1, ymax=1, #1}%
    \draw[->, figlib/thin, figlib/tip, draw=uciNavy]
      (\pgfkeysvalueof{/tikz/axes opts/xmin},0) --
      (\pgfkeysvalueof{/tikz/axes opts/xmax},0)
      node[right, figlib/font]{\pgfkeysvalueof{/tikz/axes opts/xlabel}};
    \draw[->, figlib/thin, figlib/tip, draw=uciNavy]
      (0,\pgfkeysvalueof{/tikz/axes opts/ymin}) --
      (0,\pgfkeysvalueof{/tikz/axes opts/ymax})
      node[above, figlib/font]{\pgfkeysvalueof{/tikz/axes opts/ylabel}};
  }},
  axes opts/.cd,
    xlabel/.initial=$x$,
    ylabel/.initial=$y$,
    xmin/.initial=-1, xmax/.initial=1,
    ymin/.initial=-1, ymax/.initial=1,
}

\endinput

%   \figstyle{paper}      % or \figstyle{slides}
%
% Rules:
%   • \input in preamble only — never inside tikzpicture / \fitwidth / \resizebox.
%   • Caller positions every primitive; primitives never pick their own coords.
%   • Style via figlib/* keys; never hardcode thick / \scriptsize at call sites.
%
% To extend: copy a primitive's section, document its bounding box, use figlib/*
% styles. Atoms = simple \newcommand. Pics = TikZ pic with keyval options.

\RequirePackage{etoolbox}
\usetikzlibrary{arrows.meta, calc, positioning, shapes.geometric}

% ─── Palette (idempotent: only defines if undefined) ────────────────────
\providecolor{uciGold}{HTML}{F0C040}
\providecolor{uciGoldPale}{HTML}{FAF0CC}
\providecolor{uciBlue}{HTML}{0064A4}
\providecolor{uciNavy}{HTML}{1B3D6F}
\providecolor{uciCream}{HTML}{FFF8F0}
\providecolor{uciAmber}{HTML}{C8830A}
\providecolor{uciSand}{HTML}{EED8A0}

% ─── Style switch ────────────────────────────────────────────────────────
\tikzset{
  figlib/line/.style={thick},
  figlib/thin/.style={semithick},
  figlib/font/.style={font=\scriptsize},
  figlib/tip/.style={>={Stealth[length=2mm]}},
}

\newcommand{\figstyle}[1]{%
  \ifstrequal{#1}{slides}{%
    \tikzset{%
      figlib/line/.style={thick},
      figlib/thin/.style={semithick},
      figlib/font/.style={font=\scriptsize},
      figlib/tip/.style={>={Stealth[length=2mm]}},
    }%
  }{%
    \tikzset{%
      figlib/line/.style={semithick},
      figlib/thin/.style={thin},
      figlib/font/.style={font=\footnotesize},
      figlib/tip/.style={>={Stealth[length=1.5mm]}},
    }%
  }%
}

% ============================================================
% Atoms — simple \newcommand, color via single optional arg.
% ============================================================

% \phasor[<color>]{<angle-deg>}{<magnitude>}
% Bbox: from (0,0) extending {magnitude} units at {angle}.
\newcommand{\phasor}[3][uciBlue]{%
  \draw[->, figlib/line, figlib/tip, draw=#1] (0,0) -- (#2:#3);
}

% \phasorL[<color>]{<angle-deg>}{<magnitude>}{<label>}
\newcommand{\phasorL}[4][uciBlue]{%
  \draw[->, figlib/line, figlib/tip, draw=#1] (0,0) -- (#2:#3);
  \node[figlib/font, color=#1] at (#2:{#3+0.28}) {#4};
}

% \phaseArc[<color>]{<from-deg>}{<to-deg>}{<radius>}{<label>}
% Bbox: arc of {radius} between the two angles.
\newcommand{\phaseArc}[5][uciBlue]{%
  \draw[figlib/thin, draw=#1]
    (#2:#4) arc[start angle=#2, end angle=#3, radius=#4];
  \node[figlib/font, color=#1]
    at ({(#2+#3)/2}:{#4*1.18}) {#5};
}

% \signal[<color>]{<from-coord>}{<to-coord>}
% from/to include parens, e.g. (A.east) or (1,2).
\newcommand{\signal}[3][uciNavy]{%
  \draw[->, figlib/line, figlib/tip, draw=#1] #2 -- #3;
}

% \signalL[<color>]{<from>}{<to>}{<label>}
\newcommand{\signalL}[4][uciNavy]{%
  \draw[->, figlib/line, figlib/tip, draw=#1] #2 -- #3
    node[midway, above, figlib/font, color=#1]{#4};
}

% ============================================================
% Pics — keyval options. All centered at the placement coord.
% ============================================================

% pic{block={label=..., width=..., height=...}}
% Bbox: width × height (cm, bare numbers). Default 1.0 × 0.7.
\tikzset{
  pics/block/.style={code={%
    \tikzset{block opts/.cd, label=, width=1.0, height=0.7, #1}%
    \pgfmathsetmacro{\figlibBw}{
      \pgfkeysvalueof{/tikz/block opts/width}/2}%
    \pgfmathsetmacro{\figlibBh}{
      \pgfkeysvalueof{/tikz/block opts/height}/2}%
    \draw[draw=uciNavy, fill=white, figlib/line]
      (-\figlibBw,-\figlibBh) rectangle (\figlibBw,\figlibBh);
    \node[figlib/font, color=uciNavy] at (0,0)
      {\pgfkeysvalueof{/tikz/block opts/label}};
  }},
  block opts/.cd,
    label/.initial=,
    width/.initial=1.0,
    height/.initial=0.7,
}

% pic{sum node}     — '+' in 0.5cm circle.
\tikzset{
  pics/sum node/.style={code={%
    \draw[draw=uciNavy, fill=white, figlib/line]
      (0,0) circle[radius=0.25cm];
    \node[figlib/font, color=uciNavy] at (0,0) {$+$};
  }},
}

% pic{mixer}        — '×' in 0.5cm circle.
\tikzset{
  pics/mixer/.style={code={%
    \draw[draw=uciNavy, fill=white, figlib/line]
      (0,0) circle[radius=0.25cm];
    \node[figlib/font, color=uciNavy] at (0,0) {$\times$};
  }},
}

% pic{branch node}  — 3pt filled dot.
\tikzset{
  pics/branch node/.style={code={%
    \fill[uciNavy] (0,0) circle[radius=1.5pt];
  }},
}

% pic{delay block={cycles=N, width=..., height=..., label=...}}
% N=0 → straight wire; N≥1 → that many sine cycles inside the box.
% Wire enters top mid-x, exits bottom mid-x. Label drawn LEFT of box.
% Bbox: width × height (cm, bare numbers). Default 0.55 × 0.9.
\tikzset{
  pics/delay block/.style={code={%
    \tikzset{delayblk opts/.cd,
      cycles=0, width=0.55, height=0.9, label=, #1}%
    \pgfmathsetmacro{\figlibDhalf}{
      \pgfkeysvalueof{/tikz/delayblk opts/height}/2}%
    \pgfmathsetmacro{\figlibDhalfW}{
      \pgfkeysvalueof{/tikz/delayblk opts/width}/2}%
    \pgfmathtruncatemacro{\figlibDcyc}{
      \pgfkeysvalueof{/tikz/delayblk opts/cycles}}%
    \draw[draw=uciNavy, fill=white, figlib/line]
      (-\figlibDhalfW,-\figlibDhalf)
      rectangle (\figlibDhalfW,\figlibDhalf);
    \node[figlib/font, color=uciNavy, anchor=east, inner sep=2pt]
      at (-\figlibDhalfW - 0.05, 0)
      {\pgfkeysvalueof{/tikz/delayblk opts/label}};
    \ifnum\figlibDcyc=0
      \draw[figlib/line, draw=uciBlue]
        (0,\figlibDhalf) -- (0,-\figlibDhalf);
    \else
      \pgfmathsetmacro{\figlibDmar}{0.05*
        \pgfkeysvalueof{/tikz/delayblk opts/height}}%
      \pgfmathsetmacro{\figlibDspan}{
        \pgfkeysvalueof{/tikz/delayblk opts/height}-2*\figlibDmar}%
      \pgfmathsetmacro{\figlibDstep}{\figlibDspan/(4*\figlibDcyc)}%
      \pgfmathsetmacro{\figlibDtop}{\figlibDhalf-\figlibDmar}%
      \draw[figlib/line, draw=uciBlue]
        (0,\figlibDhalf) -- (0,\figlibDtop)
        \foreach \figlibDk in {1,...,\figlibDcyc}{
          \pgfextra{%
            \pgfmathsetmacro{\figlibDya}{
              \figlibDtop-(4*\figlibDk-3)*\figlibDstep}%
            \pgfmathsetmacro{\figlibDyb}{
              \figlibDtop-(4*\figlibDk-2)*\figlibDstep}%
            \pgfmathsetmacro{\figlibDyc}{
              \figlibDtop-(4*\figlibDk-1)*\figlibDstep}%
            \pgfmathsetmacro{\figlibDyd}{
              \figlibDtop-(4*\figlibDk)*\figlibDstep}%
          }
          sin (0.13,\figlibDya) cos (0,\figlibDyb)
          sin (-0.13,\figlibDyc) cos (0,\figlibDyd)
        }
        -- (0,-\figlibDhalf);
    \fi
  }},
  delayblk opts/.cd,
    cycles/.initial=0,
    width/.initial=0.55,
    height/.initial=0.9,
    label/.initial=,
}

% pic{antenna array={N=4, spacing=1.0, labels=true, label prefix=ant\space}}
% First triangle at (0,0); subsequent at (k*spacing, 0). spacing in cm (bare).
% Bbox: x ∈ [−0.18, (N−1)*spacing+0.18], y ∈ [-0.4, 0.32] (incl labels).
\tikzset{
  pics/antenna array/.style={code={%
    \tikzset{antarr opts/.cd,
      N=4, spacing=1.0, labels=true, label prefix=ant\space, #1}%
    \pgfmathtruncatemacro{\figlibAN}{
      \pgfkeysvalueof{/tikz/antarr opts/N}-1}%
    \foreach \figlibAi in {0,...,\figlibAN}{%
      \pgfmathsetmacro{\figlibAx}{
        \figlibAi*\pgfkeysvalueof{/tikz/antarr opts/spacing}}%
      \draw[draw=uciNavy, fill=uciSand, figlib/line]
        (\figlibAx,0) -- (\figlibAx-0.18,0.32)
        -- (\figlibAx+0.18,0.32) -- cycle;
      \ifdefstring{\pgfkeysvalueof{/tikz/antarr opts/labels}}{true}{%
        \node[figlib/font, color=uciNavy, anchor=north]
          at (\figlibAx,-0.05)
          {\pgfkeysvalueof{/tikz/antarr opts/label prefix}\figlibAi};%
      }{}%
    }%
  }},
  antarr opts/.cd,
    N/.initial=4,
    spacing/.initial=1.0,
    labels/.initial=true,
    label prefix/.initial=ant\space,
}

% pic{wavefronts={N=3, length=2.4, angle=45, sep=0.7, arrow=true}}
% N parallel lines at {angle} from +x, length L, separated by sep along normal.
% First line passes through (0,0); subsequent shifted by sep along the normal.
% Optional propagation arrow drawn from (0,0) outward along -angle.
\tikzset{
  pics/wavefronts/.style={code={%
    \tikzset{wf opts/.cd,
      N=3, length=2.4, angle=45, sep=0.7, arrow=true, #1}%
    \pgfmathtruncatemacro{\figlibWN}{
      \pgfkeysvalueof{/tikz/wf opts/N}-1}%
    \pgfmathsetmacro{\figlibWnx}{
      cos(\pgfkeysvalueof{/tikz/wf opts/angle}+90)}%
    \pgfmathsetmacro{\figlibWny}{
      sin(\pgfkeysvalueof{/tikz/wf opts/angle}+90)}%
    \foreach \figlibWi in {0,...,\figlibWN}{%
      \pgfmathsetmacro{\figlibWdx}{
        \figlibWi*\pgfkeysvalueof{/tikz/wf opts/sep}*\figlibWnx}%
      \pgfmathsetmacro{\figlibWdy}{
        \figlibWi*\pgfkeysvalueof{/tikz/wf opts/sep}*\figlibWny}%
      \draw[figlib/line, draw=uciBlue!70]
        (\figlibWdx,\figlibWdy) --
          ++(\pgfkeysvalueof{/tikz/wf opts/angle}:%
             \pgfkeysvalueof{/tikz/wf opts/length});
    }%
    \ifdefstring{\pgfkeysvalueof{/tikz/wf opts/arrow}}{true}{%
      \pgfmathsetmacro{\figlibWpx}{
        0.5*\pgfkeysvalueof{/tikz/wf opts/length}*
          cos(\pgfkeysvalueof{/tikz/wf opts/angle})}%
      \pgfmathsetmacro{\figlibWpy}{
        0.5*\pgfkeysvalueof{/tikz/wf opts/length}*
          sin(\pgfkeysvalueof{/tikz/wf opts/angle})}%
      \draw[->, figlib/line, figlib/tip, draw=uciBlue]
        ($(\figlibWpx,\figlibWpy)+%
          (\pgfkeysvalueof{/tikz/wf opts/angle}+90:0.7)$)
          -- (\figlibWpx,\figlibWpy);
    }{}%
  }},
  wf opts/.cd,
    N/.initial=3,
    length/.initial=2.4,
    angle/.initial=45,
    sep/.initial=0.7,
    arrow/.initial=true,
}

% pic{beam lobe={angle=90, scale=1.0, sidelobes=true, color=uciAmber}}
% Lobe base at (0,0), main lobe pointing along {angle} from +x.
% Bbox: roughly 2*scale long, 1*scale wide.
\tikzset{
  pics/beam lobe/.style={code={%
    \tikzset{beam opts/.cd,
      angle=90, scale=1.0, sidelobes=true, color=uciAmber, #1}%
    \begin{scope}[shift={(0,0)},
                  rotate=\pgfkeysvalueof{/tikz/beam opts/angle}-90,
                  scale=\pgfkeysvalueof{/tikz/beam opts/scale}]
      \fill[\pgfkeysvalueof{/tikz/beam opts/color}!30,
            draw=\pgfkeysvalueof{/tikz/beam opts/color}, figlib/line]
        (0,0) .. controls (0.9,2.4) and (-0.9,2.4) .. (0,0);
      \ifdefstring{\pgfkeysvalueof{/tikz/beam opts/sidelobes}}{true}{%
        \fill[\pgfkeysvalueof{/tikz/beam opts/color}!30,
              draw=\pgfkeysvalueof{/tikz/beam opts/color}, figlib/line]
          (0,0) .. controls (0.9,0.7) and (0.35,0.7) .. (0,0);
        \fill[\pgfkeysvalueof{/tikz/beam opts/color}!30,
              draw=\pgfkeysvalueof{/tikz/beam opts/color}, figlib/line]
          (0,0) .. controls (-0.9,0.7) and (-0.35,0.7) .. (0,0);
      }{}%
    \end{scope}
  }},
  beam opts/.cd,
    angle/.initial=90,
    scale/.initial=1.0,
    sidelobes/.initial=true,
    color/.initial=uciAmber,
}

% pic{axes xy={xmin=-1, xmax=1, ymin=-1, ymax=1, xlabel=$x$, ylabel=$y$}}
% Origin at placement coord. Bbox: (xmax-xmin) × (ymax-ymin).
\tikzset{
  pics/axes xy/.style={code={%
    \tikzset{axes opts/.cd,
      xlabel=$x$, ylabel=$y$,
      xmin=-1, xmax=1, ymin=-1, ymax=1, #1}%
    \draw[->, figlib/thin, figlib/tip, draw=uciNavy]
      (\pgfkeysvalueof{/tikz/axes opts/xmin},0) --
      (\pgfkeysvalueof{/tikz/axes opts/xmax},0)
      node[right, figlib/font]{\pgfkeysvalueof{/tikz/axes opts/xlabel}};
    \draw[->, figlib/thin, figlib/tip, draw=uciNavy]
      (0,\pgfkeysvalueof{/tikz/axes opts/ymin}) --
      (0,\pgfkeysvalueof{/tikz/axes opts/ymax})
      node[above, figlib/font]{\pgfkeysvalueof{/tikz/axes opts/ylabel}};
  }},
  axes opts/.cd,
    xlabel/.initial=$x$,
    ylabel/.initial=$y$,
    xmin/.initial=-1, xmax/.initial=1,
    ymin/.initial=-1, ymax/.initial=1,
}

\endinput
